\section{辛形式与达布定理}
\label{sec:symplectic-forms-darboux}

现在可以引入辛形式与辛流形的一般定义了。设 $M$ 是光滑流形。

\subsection{辛形式的定义}

\begin{definition}
\label{def:symplectic-form}
\stfhIndex{I277}{xin^xing^shi}{辛形式（symplectic form）}
\stfhIndex{I281}{xin^liu^xing}{辛流形（symplectic manifold）}
$M$ 上的\emph{辛形式}是定义在 $M$ 上的一个非退化闭二形式，记为 $\omega$。二元组 $(M,\omega)$ 称为\emph{辛流形}。
\end{definition}

\begin{definition}
\label{def:symplectic-map}
\stfhIndex{I274}{xin^wei^fen^tong^pei}{辛微分同胚（symplectic diffeomorphism）}
设 $(M_1,\omega_1)$ 和 $(M_2,\omega_2)$ 是辛流形。若光滑映射（更一般地，$C^1$ 映射）$\psi:M_1\to M_2$ 满足
$\psi^*\omega_2=\omega_1$，则称 $\psi$ 是\emph{辛的}。当 $\psi$ 是微分同胚时，称它为\emph{辛微分同胚}（英文亦简称 symplectomorphism）。
\end{definition}

注意，任何辛映射自动都是\emph{浸入}。正如微分拓扑研究流形在微分同胚意义下的不变量，\emph{辛拓扑}研究辛流形在辛微分同胚意义下的不变量。关于辛几何，人们最早知道的事实之一是：与拓扑情形一样，在局部，维数是辛流形唯一的不变量。这就是著名的\emph{达布定理}。我们将在本节稍后证明它。
\stfhIndex{I092}{da^bu^ding^li}{达布定理（Darboux’ theorem）}

在继续之前，先给出一些例子。

\begin{example}
\label{ex:basic-symplectic-manifolds}
\begin{enumerate}[label=(\arabic*)]
  \item 我们已经证明，任意光滑流形 $N$ 的余切丛 $T^*N$ 都带有典范辛形式 $\omega_0$。经典相空间 $\mathbb R^n\times(\mathbb R^n)^*$ 是它的一个特例。
  \item 考虑任意可定向二维曲面 $\Sigma$ 及其一个面积形式 $\omega$，即处处非零的二形式。回忆，可定向性保证了这种二形式的存在。于是 $\omega$ 是辛形式。
  \item 复射影空间，以及更一般地复几何中的任意凯勒流形，都带有一个典范辛形式，也就是给定的凯勒形式。
\end{enumerate}
\end{example}

下面罗伯特·戈姆普夫（Robert E. Gompf）的定理特别说明，容许辛结构的四维流形构成了相当大的类。

\begin{theorem}[戈姆普夫 (Gom95)]
\label{thm:gompf-fundamental-group}
任意有限表示群都可以实现为某个辛四维流形的基本群。
\end{theorem}

现在分析辛形式定义中每项条件的后果。

首先考虑 $\omega$ 是微分二形式，即\emph{斜对称}协变二阶张量这一事实的后果。此时我们尚不假设 $\omega$ 非退化或闭。这样的二形式诱导丛映射
\[
  \widetilde\omega:TM\longrightarrow T^*M;\qquad
  X\longmapsto\omega(X,\mathord\cdot).
\]
这个映射是斜对称的。回忆，在有限维向量空间 $V$ 上有典范同构 $(V^*)^*=V$，因而伴随映射
\[
  A^*\in\operatorname{Hom}((V^*)^*,V^*)
  \cong\operatorname{Hom}(V,V^*)
\]
可以同 $A\in\operatorname{Hom}(V,V^*)$ 比较。这引出下面的一般定义。

\begin{definition}
\label{def:skew-symmetric-map}
若线性映射 $A:V\to V^*$ 满足 $A=-A^*$，则称 $A$ 是\emph{斜对称的}。
\stfhIndex{I251}{xie^dui^cheng}{斜对称（skew-symmetric）}
\end{definition}

$\omega$ 的非退化性于是蕴含 $\widetilde\omega:TM\to T^*M$ 是斜对称同构。回忆任意斜对称二次型都有偶数秩，$\omega$ 的非退化性迫使 $M$ 为偶数维。换言之，有如下结论。

\begin{corollary}
\label{cor:symplectic-even-oriented}
\stfhIndex{I196}{liu^wei^er^ce^du}{刘维尔测度（Liouville measure）}
\stfhIndex{I198}{liu^wei^er^ti^ji^xing^shi}{刘维尔体积形式（Liouville volume form）}
若 $M$ 带有辛形式，则 $M$ 必为偶数维，并带有典范体积形式 $\frac1{n!}\omega^n$。特别地，$(M,\omega)$ 可定向且具有典范定向。
\end{corollary}

我们称上述体积形式为\emph{刘维尔体积形式}，称相应的测度为\emph{刘维尔测度}。若用映射
\[
  z_j=q^j+\sqrt{-1}\,p_j,\qquad j=1,\ldots,n,
\]
来认同 $\mathbb R^{2n}\cong\mathbb R^n\times(\mathbb R^n)^*$ 与 $\mathbb C^n$，则 $\mathbb R^{2n}$ 上由刘维尔体积形式诱导的定向恰与 $\mathbb C^n$ 的复定向相同。

非退化性还使我们能够作如下定义。

\begin{definition}
\label{def:quasi-hamiltonian-poisson}
设 $\omega$ 是 $M$ 上任意一个非退化二形式，不必是闭的。取 $h\in C^\infty(M)$。
\begin{enumerate}[label=(\arabic*)]
  \item 与 $h$ 相伴的\emph{拟哈密顿向量场}记为 $X_h$，定义为
  \begin{equation}
    X_h=\widetilde\omega^{-1}(dh).
    \label{eq:2.2.4}
  \end{equation}
  \item \emph{拟泊松括号}记为 $\{f,h\}$，定义为
  \begin{equation}
    \{f,h\}=\omega(X_f,X_h).
    \label{eq:2.2.5}
  \end{equation}
\end{enumerate}
\end{definition}

先指出，等价地可以写成
\begin{equation}
  \{f,h\}=df(X_h)=X_h[f].
  \label{eq:2.2.6}
\end{equation}
从定义立刻可知，由任意非退化二形式给出的拟泊松括号都满足斜对称性、双线性与莱布尼茨法则。剩下的问题是：$\omega$ 满足什么条件时，相伴的拟泊松括号才满足雅可比恒等式？下面的结论是非退化性的一个后果。

\begin{exercise}
\label{ex:quasi-hamiltonian-ample}
证明拟哈密顿向量场的集合是充分丰富的，即对任意给定点 $x\in M$，都有
\[
  \{X(x)\mid X=X_h,\ h\in C^\infty(M)\}=T_xM.
\]
\end{exercise}

下面给出闭性在\emph{动力学}方面的第一个后果。

\begin{theorem}
\label{thm:closed-iff-jacobi}
设 $\omega$ 是非退化二形式，$\{\mathord\cdot,\mathord\cdot\}$ 是与之相伴的拟泊松括号。则 $\omega$ 是闭的，当且仅当该拟泊松括号满足雅可比恒等式。
\end{theorem}

\begin{proof}
由外微分的定义，
\begin{align}
  d\omega(X,Y,Z)
  ={}&X[\omega(Y,Z)]-Y[\omega(X,Z)]+Z[\omega(X,Y)] \notag\\
    &-\omega([X,Y],Z)+\omega([X,Z],Y)-\omega([Y,Z],X).
  \label{eq:2.2.7}
\end{align}
我们先对任意非退化闭二形式导出下面的一般恒等式。

\begin{proposition}
\label{prop:domega-jacobiator}
设 $\omega$ 如定理 \ref{thm:closed-iff-jacobi} 中所述，则
\begin{equation}
  d\omega(X_f,X_g,X_h)
  =-\bigl(\{\{g,h\},f\}+\{\{h,f\},g\}+\{\{f,g\},h\}\bigr).
  \label{eq:2.2.8}
\end{equation}
\end{proposition}

\begin{proof}
代入 $X=X_f$、$Y=X_g$、$Z=X_h$，得到
\begin{align*}
 d\omega(X_f,X_g,X_h)
 ={}&X_f[\omega(X_g,X_h)]-X_g[\omega(X_f,X_h)]+X_h[\omega(X_f,X_g)]\\
 &-\omega([X_f,X_g],X_h)+\omega([X_f,X_h],X_g)-\omega([X_g,X_h],X_f).
\end{align*}
由括号的定义，第一行变为
\[
  \{\{g,h\},f\}+\{\{h,f\},g\}+\{\{f,g\},h\}.
\]
另一方面，
\begin{align*}
 -\omega([X_f,X_g],X_h)
 &=dh([X_f,X_g])=[X_f,X_g](h)\\
 &=(\mathcal L_{X_f}\mathcal L_{X_g}-\mathcal L_{X_g}\mathcal L_{X_f})(h)\\
 &=\{\{h,g\},f\}-\{\{h,f\},g\},
\end{align*}
类似地，
\begin{align*}
 \omega([X_f,X_h],X_g)&=-\{\{g,h\},f\}+\{\{g,f\},h\},\\
 -\omega([X_g,X_h],X_f)&=\{\{f,h\},g\}-\{\{f,g\},h\}.
\end{align*}
把这些等式相加便得到 \eqref{eq:2.2.8}。
\end{proof}

由此命题立刻可知，$\omega$ 的闭性蕴含雅可比恒等式。

反过来，\eqref{eq:2.2.8} 连同拟哈密顿向量场集合的充分丰富性（习题 2.2.8）也给出逆命题；这里把细节写出。由 \eqref{eq:2.2.8}，对所有 $f,g,h$ 有
\[
  d\omega(X_f,X_g,X_h)
  =-\{\{f,g\},h\}-\{\{h,f\},g\}-\{\{g,h\},f\}=0.
\]
在任意点 $x\in M$，给定 $u,v,w$，选取满足
$X_f(x)=u$、$X_g(x)=v$、$X_h(x)=w$ 的三个拟哈密顿向量场，并在其上计算 $d\omega(u,v,w)$。这就证明雅可比恒等式蕴含 $\omega$ 的闭性，从而完成定理的证明。
\end{proof}

现在总结与辛结构 $\omega$ 相伴的泊松括号的基本性质。首先给出流形上泊松结构的一般定义。

\begin{definition}[泊松括号]
\stfhIndex{I226}{po^song^kuo^hao}{泊松括号（Poisson bracket）}
\stfhIndex{I177}{ya^ke^bi^heng^deng^shi}{雅可比恒等式（Jacobi identity）}
\label{def:poisson-bracket}
设 $M$ 是流形。$M$ 上的\emph{泊松结构}是一个称为\emph{泊松括号}的双线性映射
\[
  \{\mathord\cdot,\mathord\cdot\}:C^\infty(M)\times C^\infty(M)
  \longrightarrow C^\infty(M),
\]
并满足下列性质：
\begin{enumerate}[label=(\arabic*)]
  \item （斜对称性）$\{f,h\}=-\{h,f\}$；
  \item （双线性）$\{f+g,h\}=\{f,h\}+\{g,h\}$；
  \item （莱布尼茨法则）$\{fg,h\}=f\{g,h\}+\{f,h\}g$；
  \item （雅可比恒等式）$\{f,\{g,h\}\}+\{g,\{h,f\}\}+\{h,\{f,g\}\}=0$。
\end{enumerate}
带有泊松结构的流形 $M$ 称为\emph{泊松流形}。
\end{definition}

关于泊松流形的一般理论，可参阅 (Wn83)；关于泊松结构及其形变量子化的基本\emph{形式性定理}，可参阅 (Kon03)。有了这个一般定义，就可以陈述下面的定理。

\begin{theorem}
\label{thm:symplectic-poisson-bracket}
设 $(M,\omega)$ 是辛流形。用公式
\begin{equation}
  \{f,h\}=\omega(X_f,X_h)\;(=df(X_h))
  \label{eq:2.2.9}
\end{equation}
定义双线性映射
\[
  \{\mathord\cdot,\mathord\cdot\}:C^\infty(M)\times C^\infty(M)
  \longrightarrow C^\infty(M).
\]
那么它是定义 \ref{def:poisson-bracket} 意义下的泊松括号。此外，由辛结构诱导的泊松结构还满足：

（非退化性）对所有 $g$ 都有 $\{f,g\}=0$，当且仅当存在常值函数 $C\in\mathbb R$ 使 $f\equiv C$。
\end{theorem}

\begin{proof}
由定义 \eqref{eq:2.2.9}，斜对称性、双线性和莱布尼茨法则是显然的。另一方面，我们已经在定理 \ref{thm:closed-iff-jacobi} 中证明了雅可比恒等式。\footnote{译者注：原书此处误作“定理 2.4.8”；按上下文应为定理 2.2.9。}还需证明非退化性。证明与经典情形类似，依赖拟哈密顿向量场集合的充分丰富性，故从略。
\end{proof}

等式 $\{f,h\}=df(X_h)=X_h[f]$ 表明，当把两边看作作用在 $C^\infty(M)$ 上的导子时，有
\begin{equation}
  \{\mathord\cdot,h\}=X_h.
  \label{eq:2.2.10}
\end{equation}
这里要强调的是：按本书约定，哈密顿向量场 $X_h$ 对应的是把 $h$ 放在第二个位置的括号，而不是第一个位置。

现在，作为定理 \ref{thm:closed-iff-jacobi} 的特例，我们可以完成此前推迟到本节的经典相空间泊松括号之雅可比恒等式的证明。

\begin{corollary}
\label{cor:classical-poisson-jacobi}
定义 1.4.1 给出的经典泊松括号满足雅可比恒等式。
\end{corollary}

\begin{proof}
只需注意，由公式
\[
  \sum_{i=1}^n
  \left(
    \frac{\partial f}{\partial p_i}\frac{\partial}{\partial q^i}
    -\frac{\partial f}{\partial q^i}\frac{\partial}{\partial p_i}
  \right)
\]
给出的经典哈密顿向量场，正是与 $\mathbb R^{2n}$ 上的典范辛形式
\[
  \omega_0=\sum_jdq^j\wedge dp_j
\]
相伴的 $f$ 的哈密顿向量场。由于 $\omega_0$ 是闭的，定理 \ref{thm:closed-iff-jacobi} 即完成证明。
\end{proof}

$\omega$ 的闭性还有另一个直接后果，这一次是\emph{拓扑}方面的。

\begin{proposition}
\label{prop:closed-symplectic-nonzero-class}
设 $(M,\omega)$ 是闭辛流形，则德拉姆上同调类 $[\omega]\ne0\in H^2(M,\mathbb R)$。特别地，$H^2(M,\mathbb R)\ne0$。
\end{proposition}

\begin{proof}
注意，最高次幂 $\omega^n$ 是处处非零的 $2n$ 形式，因而定义了一个体积形式。由于 $M$ 是闭的而 $\omega^n$ 是处处非零的最高次形式，我们知道
$\int_M\omega^n>0$，所以 $[\omega^n]\ne0\in H^{2n}(M,\omega)$。利用 $H^*(M,\mathbb R)$ 在楔积下构成环这一事实，得到
\[
  [\omega]^n=[\omega^n]\ne0,
\]
从而 $[\omega]\ne0$。证明完毕。
\end{proof}

这个命题特别说明，在所有 $n$ 维球面中，只有二维球面 $S^2$ 可以成为辛流形。由于 $S^2$ 可定向，它带有一个非退化且闭的面积形式，也就是辛形式。

闭性的另一个基本后果是\emph{达布定理}；它的证明将占据本节余下部分。先从一些基本线性代数开始。

\subsection{\texorpdfstring{辛线性群 $\operatorname{Sp}(2n)$}{辛线性群 Sp(2n)}}
\stfhIndex{I280}{xin^xian^xing^qun}{辛线性群（symplectic linear group）}

考虑任意辛向量空间 $(S,\Omega)$，其中 $\Omega$ 是非退化斜对称双线性形式。若线性映射 $A:S\to S$ 满足 $A^*\Omega=\Omega$，则称 $A$ 是\emph{辛的}。注意，这种映射自动可逆。

\begin{definition}
\label{def:symplectic-linear-group}
以 $\operatorname{Sp}(S,\Omega)$ 表示 $(S,\Omega)$ 的所有辛自同构组成的集合。
\end{definition}

于是 $\operatorname{Sp}(S,\Omega)$ 是 $\operatorname{GL}(S)$（$S$ 的可逆自同构群）的一个子群。下面是线性代数的一个容易推论。

\begin{lemma}
\label{lem:symplectic-basis}
$S$ 存在一组基 $\{e_1,\ldots,e_n,f_1,\ldots,f_n\}$，使得
\begin{equation}
  \Omega(e_i,e_j)=0=\Omega(f_i,f_j),
  \qquad \Omega(e_i,f_j)=\delta_{ij}.
  \label{eq:2.2.11}
\end{equation}
\end{lemma}

任何满足上述条件的基称为\emph{辛基}或\emph{达布基}。
\stfhIndex{I272}{xin^ji}{辛基（symplectic basis）}
\stfhIndex{I090}{da^bu^ji}{达布基（Darboux basis）}

这个引理说明，任意 $(S,\Omega)$ 都同构于典范辛向量空间
\[
  S_V:=V\oplus V^*,
\]
其典范辛内积 $\Omega_V$ 定义为
\begin{equation}
  \Omega_V((v_1,\alpha_1),(v_2,\alpha_2))
  =\alpha_2(v_1)-\alpha_1(v_2).
  \label{eq:2.2.12}
\end{equation}
对 $(S_V,\Omega_V)$，任取 $V$ 的一组基 $\{e_1,\ldots,e_n\}$ 及其对偶基 $\{f_1,\ldots,f_n\}$，就在 $S_V$ 上给出一组辛基。这个选择进而给出同构
\[
  (V\oplus V^*,\Omega_V)
  \cong\left(\mathbb R^{2n},\sum_{i=1}^n dx_i\wedge dy_i\right).
\]
$\operatorname{Sp}(2n;\mathbb R)$ 是 $(\mathbb R^{2n},\omega_0)$ 的自同构群。

回忆，$\mathbb R^{2n}$ 还具有另外两种基本几何结构：记作 $(\mathord\cdot,\mathord\cdot)$ 的欧氏内积，以及复乘法
$i:\mathbb R^{2n}\to\mathbb R^{2n}$；后者通过坐标关系
\[
  z_j=x_j+iy_j
\]
把 $\mathbb R^{2n}$ 与 $\mathbb C^n$ 认同。于是它们的自同构之间有如下基本恒等式：
\[
\begin{split}
 O(2n;\mathbb R)\cap\operatorname{Sp}(2n;\mathbb R)
 &\cong O(2n;\mathbb R)\cap\operatorname{GL}(n;\mathbb C)\\
 &\cong\operatorname{GL}(n;\mathbb C)\cap\operatorname{Sp}(2n;\mathbb R)
 \cong U(n).
\end{split}
\]
下面把这些同构说得更精确。

回忆，由实化映射定义的嵌入 $\operatorname{GL}(n;\mathbb C)\hookrightarrow\operatorname{GL}(2n;\mathbb R)$ 为
\begin{equation}
  X+iY\longmapsto
  \begin{pmatrix}X&-Y\\Y&X\end{pmatrix}.
  \label{eq:2.2.13}
\end{equation}
更准确地说，对给定的 $\mathbb C$-线性映射 $A\in\operatorname{GL}(n;\mathbb C)$，其实化 $A_{\mathbb R}$ 是使下图交换的 $\mathbb R$-线性映射：
\[
\begin{array}{ccc}
 \mathbb R^n\otimes\mathbb C & \xrightarrow{\ \cong\ } & \mathbb R^{2n}\\
 {\scriptstyle A}\downarrow\phantom{\scriptstyle A} &&
 \phantom{\scriptstyle A_{\mathbb R}}\downarrow{\scriptstyle A_{\mathbb R}}\\
 \mathbb R^n\otimes\mathbb C & \xrightarrow{\ \cong\ } & \mathbb R^{2n}.
\end{array}
\]
以 $K_{\mathbb R}$ 表示子集 $K\subset\operatorname{GL}(n;\mathbb C)$ 在 $\operatorname{GL}(2n;\mathbb R)$ 中的像，并称之为 $K$ 的\emph{实化}。

同构 $O(2n;\mathbb R)\cap\operatorname{Sp}(2n;\mathbb R)\cong U(n)$ 可以在下面的引理中明确看出。

\begin{lemma}
\label{lem:orthogonal-symplectic-unitary}
矩阵 $P\in\operatorname{GL}(2n;\mathbb R)$ 属于
$O(2n;\mathbb R)\cap\operatorname{Sp}(2n;\mathbb R)$，当且仅当
\[
  P=\begin{pmatrix}X&-Y\\Y&X\end{pmatrix},
\]
并且
\[
  X^tY-Y^tX=0,
  \qquad X^tX+Y^tY=I;
\]
等价地，$P=A_{\mathbb R}$，其中 $A=X+iY\in U(n)$。换言之，
\[
  O(2n;\mathbb R)\cap\operatorname{Sp}(2n;\mathbb R)=U(n)_{\mathbb R}.
\]
\end{lemma}

将 $\operatorname{GL}(n;\mathbb C)$ 与其实化 $\operatorname{GL}(2n;\mathbb R)$ 认同后，有恒等式
\[
  \operatorname{GL}(n;\mathbb C)_{\mathbb R}\cap\operatorname{Sp}(2n;\mathbb R)
  =U(n)_{\mathbb R}.
\]

最后回忆，$U(n)$ 定义为埃尔米特内积
\begin{equation}
  \langle v,w\rangle=\sum_{i=1}^n v_i\overline{w_i}
  \label{eq:2.2.14}
\end{equation}
的自同构群，其中 $v=(v_1,\ldots,v_n)\in\mathbb C^n$，$w$ 类似。注意，此埃尔米特内积对第一个变元复线性，对第二个变元反复线性。记 $z_{\mathbb R}$ 为同构
\[
 z=(z_1,\ldots,z_n)\in\mathbb C^n
 \longmapsto(x_1,\ldots,x_n,y_1,\ldots,y_n)\in\mathbb R^{2n}
\]
的像，并称 $z_{\mathbb R}$ 为复向量 $z$ 的\emph{实化}。

\begin{lemma}
\label{lem:hermitian-real-symplectic}
设 $v,w\in\mathbb C^n$，其各自的实化为 $v_{\mathbb R},w_{\mathbb R}$。则
\[
  \langle v,w\rangle
  =(v_{\mathbb R},w_{\mathbb R})-i\omega_0(v_{\mathbb R},w_{\mathbb R}).
\]
特别地，
\[
  O(2n;\mathbb R)\cap\operatorname{Sp}(2n;\mathbb R)=U(n)_{\mathbb R}.
\]
\end{lemma}

$\operatorname{Sp}(2n;\mathbb R)$ 的另一个重要结构定理是：$U(n)_{\mathbb R}$ 是极大紧子群。因此，由李群的一般理论，
$\operatorname{Sp}(2n;\mathbb R)/U(n)_{\mathbb R}$ 可缩。（例如参见 (He78)；关于从 $\operatorname{Sp}(2n;\mathbb R)$ 到 $U(n)_{\mathbb R}$ 的一个显式形变收缩的构造，可参阅 (CZ84, SZ92)。）这蕴含 $\operatorname{Sp}(2n;\mathbb R)$ 与 $U(n)$ 具有相同的同伦型。固定一个从 $\operatorname{Sp}(2n;\mathbb R)$ 到 $U(n)_{\mathbb R}$ 的形变收缩后，可以把行列式映射 $\det:U(n)\to S^1$ 延拓为
\[
  \widetilde{\det}:\operatorname{Sp}(2n;\mathbb R)\longrightarrow S^1,
\]
它在 $\pi_1$ 上诱导同构。显然，该映射连续；但与 $\det:U(n)\to S^1$ 不同，它并不是同态。（关于这个问题的相关讨论参见 (CZ84, SZ92)。）这是因为形变收缩 $\operatorname{Sp}(2n;\mathbb R)\to U(n)_{\mathbb R}$ 不保持群结构。（参见 (SZ92)。）

\begin{exercise}
\label{ex:no-group-retraction-sp-u}
证明不存在从 $\operatorname{Sp}(2n;\mathbb R)$ 到 $U(n)$ 且保持群结构的形变收缩。
\end{exercise}

\begin{definition}
\label{def:symplectic-loop-index}
设 $\gamma$ 是 $\operatorname{Sp}(2n;\mathbb R)$ 中的环路。指标 $\mu(\gamma)$ 定义为映射
\[
  \widetilde{\det}\circ\gamma:S^1\longrightarrow S^1
\]
的度数。
\end{definition}

\subsection{达布定理的证明}

现在陈述辛流形的基本达布定理。
\stfhIndex{I092}{da^bu^ding^li}{达布定理（Darboux’ theorem）}

\begin{theorem}[达布定理]
\label{thm:darboux}
设 $(M,\omega)$ 是任意 $2n$ 维辛流形。对每一点 $x_0\in M$，都存在包含 $x_0$ 的坐标邻域以及坐标卡 $\varphi:U\to\mathbb R^{2n}$，使 $\omega=\varphi^*\omega_0$；也就是说，对
$\varphi=(x_1,\ldots,x_n,y_1,\ldots,y_n)$ 有
\[
  \omega=\sum_{i=1}^n dx_i\wedge dy_i.
\]
任何这样的坐标都称为\emph{达布坐标}。
\stfhIndex{I091}{da^bu^zuo^biao}{达布坐标（Darboux coordinates）}
\end{theorem}

\begin{proof}
证明方案使用辛几何的一项基本技巧，称为\emph{莫泽形变法} (Mo65)。
\stfhIndex{I208}{mo^ze^xing^bian^fa}{莫泽形变法（Moser’s deformation method）}

第一步，选取微分同胚
\[
  \psi:U'\subset M\longrightarrow V'\subset T^*\mathbb R^n\cong\mathbb R^{2n},
\]
使得在给定点 $x_0\in U'$ 有
\begin{equation}
  \omega|_{x_0}=\psi^*\omega_0|_{x_0}.
  \label{eq:2.2.15}
\end{equation}
这是因为在合同作用下，$\operatorname{GL}(2n)$ 在 $\mathbb R^{2n}$ 上斜对称双线性形式的集合上具有传递作用。

为实现 \eqref{eq:2.2.15}，先任取局部坐标卡
$\widetilde\psi:\widetilde U\to\widetilde V$，其中 $x_0\in\widetilde U\subset U'$。于是 $T_{\widetilde\psi(x_0)}\mathbb R^{2n}\cong\mathbb R^{2n}$ 上有两个非退化二次型：一个是 $\omega_0$，另一个是 $\omega_1:=\widetilde\psi_*\omega$。任意两个非退化斜对称双线性形式彼此合同，所以存在矩阵 $A\in\operatorname{GL}(2n)$，使
\[
  A^*\omega_1:=\omega_1(A\mathord\cdot,A\mathord\cdot)=\omega_0.
\]
于是取 $\psi=A^{-1}\circ\widetilde\psi$，它便满足 \eqref{eq:2.2.15}。

第二步，把 $\psi$ 形变为所需的 $\varphi:U\to\mathbb R^{2n}$。为此，将 $\psi$ 与一个满足
\[
  h^*(\psi^*\omega_0)=\omega
\]
的微分同胚 $h:U\to U'$ 复合，其中 $U\subset U'$。于是 $\varphi=\psi\circ h$ 达到目的。

为完成这一步，我们将证明温斯坦 (Wn73) 的一个一般形变引理。它比证明达布定理实际所需的结论更强；在达布定理中，$N$ 只是单点集 $\{x_0\}$。

\begin{lemma}[形变引理]
\label{lem:weinstein-deformation}
设 $N\subset M$ 是紧子流形，$\omega$ 与 $\omega'$ 是定义在 $N$ 的某个邻域上的两个辛形式。若
$\omega\equiv\omega'$ 在 $TM|_N$ 上成立，则存在微分同胚 $h:U\to U'$，其中 $U,U'$ 都是 $N$ 的邻域，并且
\begin{enumerate}[label=(\arabic*)]
  \item $dh|_{TM|_N}\equiv\operatorname{id}|_{TM|_N}$；
  \item 在 $U\supset N$ 上有 $h^*\omega'=\omega$。
\end{enumerate}
\end{lemma}

\noindent\emph{莫泽形变法。}
考虑单参数族
\[
  \omega_t=(1-t)\omega+t\omega',\qquad 0\le t\le1.
\]
注意，对所有 $0\le t\le1$，在 $TM|_N$ 上都有
$\omega_t|_{TM|_N}\equiv\omega|_{TM|_N}(=\omega'|_{TM|_N})$。特别地，在 $N$ 的某个小邻域内，所有 $\omega_t$ 都非退化，因而都是辛形式。我们寻找定义在 $N$ 附近的一族微分同胚 $\{h_t\}_{0\le t\le1}$，使
\begin{equation}
 \left\{
 \begin{aligned}
  h_t^*\omega_t&=\omega,\\
  dh_t|_{TM|_N}&=\operatorname{id}\quad\text{于 }TM|_N.
 \end{aligned}
 \right.
 \label{eq:2.2.16}
\end{equation}
莫泽形变法 (Mo65) 的关键是：不直接求解这个问题，而把它转化为求生成同痕 $h_t$ 的向量场族
\[
  \xi_t=\frac{dh_t}{dt}\circ h_t^{-1}.
\]
当然，一旦得到 $\xi_t$，就可以通过求解常微分方程初值问题
\[
  \frac{dh_t}{dt}=\xi_t\circ h_t,
  \qquad h_0=\operatorname{id}
\]
恢复 $h_t$。为写出 $\xi_t$ 应满足的条件，对 \eqref{eq:2.2.16} 关于 $t$ 求导：
\begin{align*}
 0&=\frac d{dt}h_t^*\omega_t
   =h_t^*\left(\mathcal L_{\xi_t}\omega_t+\frac{d\omega_t}{dt}\right)\\
  &=h_t^*\left(d(\xi_t\mathbin{\lrcorner}\omega_t)+\frac{d\omega_t}{dt}\right).
\end{align*}
注意，由于 $h_t|_N\equiv\operatorname{id}|_N$ 且 $[0,1]$ 紧，对每个 $t\in[0,1]$，$h_t$ 都把 $N$ 的一个固定邻域 $U_1$ 映入 $N$ 的另一个较大邻域 $U_2$。因此可以去掉 $h_t^*$，得到
\begin{equation}
 \left\{
 \begin{aligned}
  d(\xi_t\mathbin{\lrcorner}\omega_t)+d\omega_t/dt&=0
    &&\text{于 }U_1,\\
  \xi_t&\equiv0.
 \end{aligned}
 \right.
 \label{eq:2.2.17}
\end{equation}
令 $\beta_t=\xi_t\mathbin{\lrcorner}\omega_t$，并注意 $-d\omega_t/dt=\omega-\omega'$，就把 \eqref{eq:2.2.17} 进一步化为
\begin{equation}
 \left\{
 \begin{aligned}
  d\beta_t&=\omega-\omega' &&\text{于 }N\text{ 的某邻域},\\
  \beta_t&\equiv0 &&\text{于 }N.
 \end{aligned}
 \right.
 \label{eq:2.2.18}
\end{equation}
这里，方程 $d\beta_t=\omega-\omega'$ 究竟在哪个邻域上求解，也将一并确定。作为管状邻域定理的一个简单应用，可以找到 $N$ 的一个管状邻域 $U_1\supset N$ 到 $N$ 上的形变收缩
\[
  \psi_t:U_1\longrightarrow U,
  \qquad 0\le t\le1,
  \qquad \psi_t|_N=\operatorname{id}_N,
\]
使 $t>0$ 时 $\psi_t$ 都是到其像上的微分同胚，且 $\psi_1=\operatorname{id}|_{U_1}$。

现在令 $\Omega=\omega-\omega'$，并考虑 $\psi_t^*\Omega$。于是
\[
\begin{aligned}
 \Omega=\psi_1^*\Omega
 &=\psi_0^*\Omega+\int_0^1\frac d{dt}\psi_t^*\Omega\,dt\\
 &=\psi_0^*\Omega+\int_0^1\psi_t^*(\eta_t\mathbin{\lrcorner}\Omega)\,dt,
\end{aligned}
\]
其中
\[
 \eta_t=
 \begin{cases}
  d\psi_t/dt\circ\psi_t^{-1},&t>0,\\
  0,&t=0.
 \end{cases}
\]
由假设 $\omega\equiv\omega'$ 于 $TM|_N$，以及 $\operatorname{Im}\psi_0\subset N$，可知 $\psi_0^*\Omega=0$。因此只需把形式 $\beta_t$ 取为
\[
  \beta_t\equiv\int_0^t\psi_s^*(\eta_s\mathbin{\lrcorner}\Omega)\,ds,
\]
并令其与 $t$ 无关。注意，由于 $\psi_t|_N=\operatorname{id}|_N$ 且 $\eta_t\equiv0$ 于 $N$，所以 $\beta_t=0$ 于 $N$。这样就构造了满足 \eqref{eq:2.2.18} 的 $\beta_t$，从而也构造了满足 \eqref{eq:2.2.17} 的向量场 $\xi_t$。

最后，为了在 $N$ 的一个邻域上得到同痕 $h_t$，必须保证：定义于 $[0,1]\times U_1$ 上的常微分方程 $\dot x=\xi_t(x)$ 的解之存在域
$\mathcal D\subset[0,1]\times U_1$，包含一个形如 $[0,1]\times U$ 且包含于 $\mathcal D$ 的子集。由于 $\xi_t\equiv0$ 于 $N$，且假设 $N$ 紧，$\mathcal D$ 包含 $[0,1]\times N$。另一方面，由常微分方程的基本存在定理，$\mathcal D$ 总是 $[0,1]\times U_1$ 的开子集。又因 $[0,1]$ 紧，$\mathcal D$ 必包含一个形如 $[0,1]\times U$ 的子集，其中 $N\subset U\subset U_1$。这证明了对所有 $(t,x)\in[0,1]\times U$，$h_t:U\to h_t(U)$ 都有定义，并满足所需性质 \eqref{eq:2.2.16}。

一般形变引理的证明完成后，令 $N=\{x\}\subset M$，便完成了达布定理的证明。
\end{proof}

现在回过头来陈述莫泽定理 (Mo65)；辛形式的整个同痕问题正是从这里开始的。先引入以下一般定义。

\begin{definition}
\label{def:isotopic-symplectic-forms}
\stfhIndex{I175}{tong^hen}{同痕（isotopic）}
\stfhIndex{I232}{wei^tong^hen}{伪同痕（pseudo-isotopic）}
设 $M$ 是闭流形，$\omega_0,\omega_1$ 是 $M$ 上的两个辛形式。
\begin{enumerate}[label=(\arabic*)]
  \item 若存在辛形式的光滑族 $\{\omega_t\}$，使其上同调类满足
  \[
    [\omega_t]=[\omega_0]\quad\text{于 }H^2(M,\mathbb R),
    \qquad t\in[0,1],
  \]
  则称 $\omega_0$ 与 $\omega_1$ \emph{同痕}。称这样的族 $\{\omega_t\}$ 为 $\omega_0$ 与 $\omega_1$ 之间的\emph{同痕}。
  \item 若存在连接 $\omega_0$ 与 $\omega_1$ 的辛形式光滑族，而不要求其上同调类保持不变，则称 $\omega_0$ 与 $\omega_1$ \emph{伪同痕}。
  \item 若存在微分同胚 $\varphi:M\to M$，使 $\omega_0$ 与 $\varphi^*\omega$ 彼此伪同痕，则称 $\omega_0$ 与 $\omega_1$ \emph{形变等价}。
  \stfhIndex{I098}{xing^bian^deng^jia}{形变等价（deformation equivalent）}
\end{enumerate}
\end{definition}

\begin{exercise}[莫泽]
\label{ex:moser-isotopy}
设 $M$ 是闭流形。对任意两个同痕的形式 $\omega_0,\omega_1$ 以及一个同痕 $\{\omega_t\}$，都存在微分同胚的同痕 $\varphi_t:M\to M$，满足
$\varphi_0=\operatorname{id}$ 且 $\varphi_t^*\omega_t=\omega_0$。特别地，$\omega_0$ 与 $\omega_1$ 微分同胚。（提示：使用霍奇分解定理，或参见 (Mo65)。）
\end{exercise}

达布定理有如下直接推论：$\mathbb C^n$ 中充分小的标准球可以辛嵌入 $(M,\omega)$。

\begin{corollary}
\label{cor:small-symplectic-ball}
设 $(M,\omega)$ 是任意辛流形。对充分小的 $R>0$，存在辛嵌入
\[
  h:B^{2n}(R)\subset\mathbb C^n\longrightarrow(M,\omega).
\]
\end{corollary}

这引出了 $(M,\omega)$ 的格罗莫夫半径这一概念。
\stfhIndex{I143}{ge^luo^mo^fu^ban^jing}{格罗莫夫半径（Gromov radius）}

\begin{definition}
\label{def:gromov-radius}
使得存在辛嵌入
\[
  h:B^{2n}(R)\subset\mathbb C^n\longrightarrow(M,\omega)
\]
的所有 $R$ 的上确界，称为\emph{格罗莫夫半径}，记作 $\widetilde r(M,\omega)$。
\end{definition}

在这个意义下，格罗莫夫定理 (Gr85) 是揭示辛映射与保体积映射之间本质差异的基本定理。

\begin{theorem}[格罗莫夫非挤压定理]
\stfhIndex{I145}{ge^luo^mo^fu^fei^ji^ya^ding^li}{格罗莫夫非挤压定理（Gromov’s non-squeezing theorem）}
\label{thm:gromov-nonsqueezing}
考虑带有标准辛形式 $\omega_0$ 的柱体
\[
  Z_n(r)=D^2(r)\times\mathbb C^{n-1}\subset\mathbb C^n.
\]
则
\[
  \widetilde r(Z_n(r),\omega_0)=r.
\]
\end{theorem}

其完整证明见第 11.1 节，那里沿用格罗莫夫在 (Gr85) 中通过伪全纯曲线方法给出的原始几何证明。关于通过研究哈密顿周期轨道得到的动力学证明，可参阅 (EkH89)。在此阶段，我们要强调，若能发展出一种同时容纳辛几何这两个方面的方法，那将是极其理想的。本书的主要目标之一，正是通过\emph{链层弗洛尔理论}传达这一数学观点。
